\section{Considerações Finais}
\textbf{(1) Lexical e semântico são equivalentes em média e complementares por consulta.} BM25 e bi-encoder têm NDCG@10 de \NdcgTenLex{} e \NdcgTenSem{} no \emph{test} (diferença sem significância), mas ganham em consultas diferentes: o lexical vence quando há termo raro ou nome próprio (\emph{kohlrabi}, \emph{shelf life}, \emph{grapes}), e o semântico quando há sinonímia, abreviação, relevância indireta ou palavras raras sem valor temático (\emph{low-carb}, \emph{cocaine}, \emph{vitamin D recommendation}). Os Top-10 se sobrepõem pouco (Jaccard \OverlapTop).
\textbf{(2) O RRF é a melhoria mais barata e consistente:} $+0{,}025$ de NDCG@10 sobre o lexical e $+0{,}039$ sobre o semântico, com sensibilidade pequena a $k$ e custo de décimos de milissegundo.
\textbf{(3) A correção ortográfica funciona, mas é frágil em empates:} recuperou \TypoRecAgg\% da queda de NDCG@10 nas 10 consultas do \emph{test} (95\% no \emph{dev}); a única falha foi uma troca de letras cujo candidato correto perdeu no desempate por Jaccard e frequência. Os embeddings também sofrem com erros de digitação.
\textbf{(4) O reranking com o cross-encoder MS MARCO MiniLM não compensou o custo neste experimento:} NDCG@10 de \NdcgTenHyb{} para \NdcgTenRr{} ($p=0{,}96$), P@10 menor, e \CEMsQuery{} s por consulta em CPU. Há margem (oráculo de \OracleNdcg{} no Top-50), mas o modelo a aproveita mal; vale testar um cross-encoder treinado no domínio biomédico e mais consultas do \emph{dev} na escolha.
\textbf{(5) RAG:} a recuperação entregou ao menos um relevante ao contexto em 9 de 10 consultas, e a única exceção (PLAIN-2850) é uma falha de recuperação, não de geração. \todo{Concluir, após a avaliação humana, o efeito do RAG na fidelidade, completude, utilidade e rastreabilidade.}
\textbf{Trabalhos futuros:} correção sensível ao contexto, expansão de consultas para termos fora do vocabulário (nomes de remédios, alimentos e lugares), fusão ponderada ou aprendida, reranker específico de domínio e avaliação do RAG com mais consultas e mais de uma LLM.
